\begin{lemma}[Bridge response transition]
Assume $F$ is a front on the full base and
$\mathsf{FrontBridge}(F,s,t,u)$.  Choose the next left node by
\[
 (s\in F\text{ and }s'=s)\quad\text{or}\quad
 (s\notin F\text{ and }s'=u).
\]
Suppose the right response is one of the following:
\begin{enumerate}
\item $t\notin F$, $t'$ is an ordered proper one-point extension of the
      nonempty node $t$ in $T(F)$, and for its new point $n$ we have
      $j<n$ for every $j\in u$ and set $u'=u\cup\{n\}$; or
\item $t\in F$, $t'=t$, and we choose a fresh $k$ larger than every member
      of $u$, setting $u'=u\cup\{k\}$.
\end{enumerate}
Then $\mathsf{FrontBridge}(F,s',t',u')$ holds.  The selected left node is
legal, and it is equal to $s$ when $s$ is terminal and is an exact ordered
one-point extension of $s$ otherwise.
\end{lemma}
\begin{proof}
Unpack the eight clauses of \noderef{FrontBridge}.  We first record the
finite-set calculation used below.  If $q$ is larger than every member of
the nonempty set $u$ and $u_q=\operatorname{insert}(q,u)$, then the
definition of \noderef{FrontBridgeTail} gives
\[
 \mathsf{FrontBridgeTail}(u_q)
 =\operatorname{insert}\bigl(q,\mathsf{FrontBridgeTail}(u)\bigr).
 \tag{1}
\]
Indeed, $q$ belongs to the left-hand side because $u$ is nonempty, so some
member of $u$ is below $q$.  If $j\in u$, then $j$ belongs to the new tail
exactly when some member of $u$ is below $j$: the additional possible
witness $q<j$ is impossible because $j<q$.  This is precisely (1), using
the definition in \noderef{FrontBridgeTail}.

We also spell out how an old initial segment transfers when $q$ is appended
above $u$.  If $\mathsf{InitialSegment}(v,u)$, then by
\noderef{InitialSegment} either $v=u$, or there is a cutoff $c\in u$ with
\[
 v=\{j\in u\mid j<c\}.
 \tag{2}
\]
In the equality case, every member of $u$ is below $q$, so
\[
 v=u=\{j\in u_q\mid j<q\},
\]
and $q$ is a cutoff in $u_q$.  In the second case, $c\in u_q$ and
$c<q$, and therefore
\[
 \{j\in u_q\mid j<c\}=\{j\in u\mid j<c\}=v.
 \tag{3}
\]
Thus $\mathsf{InitialSegment}(v,u_q)$ follows in both cases, with an
explicit cutoff (respectively $q$ or $c$).

Consider first the left move.  If $s\in F$, then $s'=s$ by the hypothesis
on the move, $s'\in\mathsf{PrefixTree}(F)$ by the bridge, and the bridge
gives $\mathsf{InitialSegment}(s,u)$.  Applying (2)--(3) with the new
point $q$ from the relevant right response gives
$\mathsf{InitialSegment}(s',u')$.

If $s\notin F$, the bridge supplies $a$ such that, with the predicate
$\mathsf{ProperInitialSegment}$ understood as in \noderef{ProperInitialSegment},
\[
 u=\operatorname{insert}(a,s),\qquad
 \mathsf{ProperInitialSegment}(s,\operatorname{insert}(a,s)),\qquad
 u\in\mathsf{PrefixTree}(F).
 \tag{4}
\]
Here $s'=u$ by the left-move hypothesis.  In either right response the new
point $q$ satisfies $j<q$ for every $j\in u$, so the cutoff $q$ directly
shows
\[
 \mathsf{InitialSegment}(u,\operatorname{insert}(q,u)),
 \qquad
 \mathsf{ProperInitialSegment}(u,\operatorname{insert}(q,u)).
 \tag{5}
\]
If $u\in F$, its membership in the prefix tree is immediate (and also
already supplied by (4)); if $u\notin F$, apply
\noderef{FrontTreeOrderedExtension} to the tree membership in (4) and the
proper extension in (5).  In either subcase $u'=\operatorname{insert}(q,u)$
belongs to the prefix tree.  Consequently the new left node satisfies the
left clause of \noderef{FrontBridge}, whether it is terminal (use the first
part of (5)) or nonterminal (use the second part of (5) and the displayed
insert equality).

Now take the first right response.  It supplies
\[
 t'=\operatorname{insert}(n,t),\qquad
 u'=\operatorname{insert}(n,u),\qquad
 j<n\quad(j\in u).
 \tag{6}
\]
Since $t\notin F$, the bridge gives
$t=\mathsf{FrontBridgeTail}(u)$.  Applying (1) with $q=n$ yields
\[
 \mathsf{FrontBridgeTail}(u')
 =\operatorname{insert}(n,\mathsf{FrontBridgeTail}(u))
 =\operatorname{insert}(n,t)=t'.
 \tag{7}
\]
If $t'\notin F$, (7) is exactly the required tail clause.  If $t'\in F$,
(7) implies $\mathsf{InitialSegment}(t',\mathsf{FrontBridgeTail}(u'))$
by the equality alternative in \noderef{InitialSegment}.  The response
already supplies $t'\in\mathsf{PrefixTree}(F)$, and $t'$ is nonempty
because it extends the nonempty $t$; $u'$ is nonempty because $u$ is.
Together with the left analysis above, these are all clauses of
\noderef{FrontBridge}, so the new bridge holds in this case.

Finally take the fresh-$k$ response.  Here $t\in F$, $t'=t$, and
$u'=\operatorname{insert}(k,u)$ with every member of $u$ below $k$.
The bridge gives
\mathsf{InitialSegment}(t,\mathsf{FrontBridgeTail}(u)).  Apply the
explicit cutoff analysis (2)--(3) not to $u$ but to the tail
$\mathsf{FrontBridgeTail}(u)$, whose appended point is $k$.  More
explicitly, if $t=\mathsf{FrontBridgeTail}(u)$, then (1) gives
\[
 \mathsf{FrontBridgeTail}(u')
 =\operatorname{insert}(k,t),
 \]
and $t=\{j\in\mathsf{FrontBridgeTail}(u')\mid j<k\}$, so the required
initial segment uses the new cutoff $k$.  Otherwise the definition of
\noderef{InitialSegment} supplies $c\in\mathsf{FrontBridgeTail}(u)$ with
$t=\{j\in\mathsf{FrontBridgeTail}(u)\mid j<c\}$; since $c<k$, the same
cutoff $c$ gives
\[
 t=\{j\in\mathsf{FrontBridgeTail}(u')\mid j<c\}.
\]
Thus $\mathsf{InitialSegment}(t',\mathsf{FrontBridgeTail}(u'))$ holds,
while $t'=t\in\mathsf{PrefixTree}(F)$ and $t'$ remains nonempty.  The
left analysis with $q=k$ and the nonemptiness of $u'$ finish all clauses
of \noderef{FrontBridge} in the second case as well.

It remains only to record the final assertion about the left move.  If
$s\in F$, the choice gives $s'=s$.  If $s\notin F$, use (4) and
$s'=u$ to take the same witness $a$: then
$s'=\operatorname{insert}(a,s)$ and
$\mathsf{ProperInitialSegment}(s,\operatorname{insert}(a,s))$ by
\noderef{ProperInitialSegment}.
This proves the stated legality and exact-extension conclusion.
\end{proof}
